% Versión en español del perfil DevOps Engineer. Se mantiene afinado a roles de
% infraestructura/cloud (AWS, Terraform, CI/CD). Para roles de liderazgo
% centrados en remediación de deuda técnica y gobernanza, ver lang/lead/es.tex.
%
% Mantener cifras y alcance alineados con lang/en.tex. La frontera de hechos
% confirmados / no declarados que respetan todos los perfiles está en FACTS.md.

% Versión en español. Mantener cifras y alcance alineados con lang/en.tex.
\resumeaddress{Ciudad de México, México}
\resumejobtitle{Ingeniero DevOps \textbar\ AWS, Terraform, Kubernetes}
\resumesummary{Ingeniero DevOps con trayectoria en desarrollo de software desde 2017 y experiencia en producción en banca y telecomunicaciones. Aprovisiono infraestructura en AWS y Azure como código con Terraform, Ansible y Python/Go, y administro 9 clústeres de QA y 12 clústeres de preproducción (DEV, INTGR y Certificación) sobre Kubernetes y OpenShift. Implemento y opero controles \textbf{ISO 27001} y reviso la arquitectura de microservicios del equipo. Aporto la perspectiva de un desarrollador de aplicaciones al trabajo de infraestructura, desde el código en Java/Spring Cloud hasta la operación de clústeres.}
\resumesummaryheading{Sobre mí}
\makeheader

% Sidebar column. In the ATS layout \sidecol is a no-op and this prints inline.
\sidecol{%
\section*{Competencias Técnicas}
\skillrow{Plataformas y CI/CD}{Kubernetes, OpenShift, Docker, Helm, Azure DevOps, Jenkins}
\skillrow{AWS e infraestructura}{AWS (EC2, S3, VPC, Security Groups, IAM, RDS, CloudWatch, Route~53), Terraform, Azure, Ansible, Linux, Bash, Python, Go}
\skillrow{Mensajería}{Kafka, RabbitMQ, IBM~MQ}
\skillrow{Datos}{Oracle, PostgreSQL}
\skillrow{Observabilidad, seguridad y cumplimiento}{Elastic Stack (ELK), Splunk, ServiceNow, ISO 27001 (implementación y operación de controles), SonarQube Cloud, Snyk, Mend}
\skillrow{Desarrollo}{Java, Spring Cloud, JPA, APIs REST, TypeScript, React, Node.js}
\skillrow{Entrega y prácticas}{Scrum, Microsoft Project, Git, Maven, Nexus, GitHub}
}

% Main column.
\maincol{%
\section*{Experiencia Profesional}

\resumeentry{Ingeniero DevOps}{Veradat Smart Network S.A. de C.V.}{Ago 2024 -- Oct 2026}{Ciudad de México}
\begin{itemize}
    \item Diseñé y mantuve \textbf{10 pipelines de CI/CD en Azure DevOps} para \textbf{3 equipos} (frontend, backend, QA), reduciendo el tiempo de despliegue del conjunto completo de \textbf{70 microservicios} de \textbf{6 horas a 1 hora} (aproximadamente 83\% menos tiempo). Desarrollé los \textbf{charts de Helm} que los tres equipos adoptaron como estándar de despliegue, y registré el trabajo de infraestructura en \textbf{Microsoft Project} por sprint.
    \item Revisé la \textbf{arquitectura hexagonal} del equipo, que desacopla cada microservicio de sus dependencias tecnológicas (PostgreSQL, Oracle, RabbitMQ, Kafka, IBM~MQ, SQS), y generé la \textbf{imagen Docker para cada combinación tecnológica} que cada microservicio requiere.
    \item Administro \textbf{9 clústeres de QA: 8 Kubernetes y 1 OpenShift}, con sondas readiness/liveness que controlan cuándo las instancias reciben tráfico y disparan su recuperación en actualizaciones progresivas, y monitoreo y alertas con \textbf{Elastic Stack} que redujeron el \textbf{tiempo de detección de incidentes de QA de 60 a 30 minutos} (50\% menos tiempo).
    \item Aprovisioné infraestructura en \textbf{AWS y Azure con Terraform} -- \textbf{EC2, S3, VPC, Security Groups, IAM, RDS, CloudWatch y Route~53} en AWS, junto con máquinas virtuales, balanceadores de carga y gestión de secretos en ambas nubes -- y automaticé la configuración de 10 servidores Linux con Ansible, reduciendo el aprovisionamiento de \textbf{2 horas a 20 minutos} (aproximadamente 83\% menos tiempo). Operé \textbf{4 clústeres de DEV, 4 de INTGR y 4 de Certificación (preproducción)} junto a los 9 de QA.
    \item Realicé \textbf{backups y rollbacks de bases de datos PostgreSQL y Oracle} que respaldan los 70 microservicios, asegurando la recuperabilidad de los datos durante despliegues y respuesta a incidentes.
    \item \textbf{Implementé y operé controles ISO 27001} en gestión de secretos, control de acceso en nube y escaneo de dependencias, y remedié vulnerabilidades críticas y altas detectadas por SonarQube Cloud, Snyk y Mend mediante actualizaciones de paquetes del sistema operativo y dependencias Maven.
    \item Desarrollé \textbf{scripts de automatización e integración en Python y Go} para apoyar operaciones de infraestructura, monitoreo y herramientas internas, y utilicé \textbf{Claude Code} como asistente de desarrollo.
    \item \textbf{Coordiné la interlocución técnica} con stakeholders del cliente, incluido el reporte formal de seguridad que recibe su área de seguridad, y entre los \textbf{3 equipos} (frontend, backend, QA), mediante \textbf{dashboards y boards de Azure DevOps} y una \textbf{wiki compartida} como canal asincrónico.
\end{itemize}

\resumeentry{Ingeniero Fullstack + DevOps}{IDS Comercial S.A. de C.V. \textbar\ Cliente: Banamex}{Ene 2021 -- Ago 2024}{México, remoto}
\begin{itemize}
    \item Desarrollé y mantuve \textbf{10 microservicios Java/Spring Cloud} para flujos de trabajo de ejecutivos bancarios. Esos flujos atendían aproximadamente \textbf{1,000 transacciones diarias por sucursal}.
    \item \textbf{Estandaricé los microservicios con un POM padre multimódulo de Maven} y gestión de dependencias compartida, garantizando un versionado consistente y reduciendo el esfuerzo de actualización en todos los servicios.
    \item \textbf{Me apegaba a los controles ISO 27001} en este entorno bancario durante toda mi colaboración.
    \item Construí pipelines de CI/CD con \textbf{Jenkins} para automatizar la compilación y el despliegue de aplicaciones bancarias. Desarrollé funcionalidades frontend con React, Node.js y TypeScript, y administré bases de datos Oracle.
    \item Atendí \textbf{tickets de infraestructura en ServiceNow} bajo un \textbf{SLA de 1 día}, revisando logs para \textbf{troubleshooting} y \textbf{levantando incidentes al área correspondiente}. Monitoreé producción con Splunk y atendí incidentes del entorno de QA (tiempo promedio de resolución de 2 a 3 horas; casos complejos de hasta una semana).
\end{itemize}

\resumeentry{Ingeniero de Software Fullstack}{Proyecto continuo para Telcel, mediante GLOBAL HITSS y Everis}{Jul 2018 -- Dic 2020}{Ciudad de México}
\textit{GLOBAL HITSS: jul 2018 -- sep 2019; Everis: sep 2019 -- dic 2020.}\par\vspace{3pt}
\begin{itemize}
    \item Migré \textbf{microservicios de un servidor de aplicaciones Tomcat a OpenShift}, conteniendo la capa de integración construida con \textbf{Java, Spring Cloud y JPA} que traducía protocolos de sistemas legados de telecomunicaciones a APIs REST para conectarlos con servicios modernos.
    \item Implementé pipelines de CI/CD y contenericé el servicio con \textbf{Docker sobre Kubernetes}. Automaticé procesos operativos con Bash y gestioné repositorios de código y artefactos con GitHub y Nexus.
\end{itemize}

\section*{Proyecto Destacado}
\resumeentry{Sistema de Gestión de Cumplimiento de la Ley de Subcontratación}{GDP administración S de RL de CV (Grupo de la Parra)}{Jun 2021 -- Dic 2021}{Ciudad de México}
\begin{itemize}
    \item Lideré, como líder técnico, el diseño técnico y la arquitectura end-to-end de una plataforma en producción para gestionar el cumplimiento de la \textbf{Ley Federal del Trabajo en materia de subcontratación}, cubriendo una fuerza laboral subcontratada de \textbf{5,000 empleados} y múltiples empresas proveedoras y receptoras de servicios especializados.
    \item Diseñé un \textbf{motor de auditoría de cumplimiento} que valida reglas regulatorias (existencia de CFDI por empleado vs. fechas de pago, pagos de IMSS/INFONAVIT/retenciones vs. calendario, alta ante la Secretaría del Trabajo) y emite \textbf{alertas automatizadas} por vencimientos de contratos y plazos de entrega de información.
    \item Implementé un \textbf{panel de BI de cumplimiento} con reportes regulatorios (\textbf{SAT, IMSS, INFONAVIT, STPS}) y un dashboard con KPIs (documentos pendientes, días para entrega, contratos por proveedor/cliente) que permitió al área legal/antifraude operar de forma proactiva.
\end{itemize}
}

% Back to the sidebar.
\sidecol{%
\section*{Educación}
\resumeentry{IPN -- Escuela Superior de Cómputo (ESCOM)}{Ingeniería en Sistemas Computacionales}{2014 -- 2018}{Ciudad de México}

\section*{Certificación, Idiomas y Programación}
\skillrow{Certificación}{Java SE -- Oracle, Certificado Asociado (2013). \href{https://www.youracclaim.com/badges/583548b1-06fa-4dbf-889d-419e6c6c2ec1/public_url}{Credencial}; Scrum Fundamentals -- SCRUMstudy (2018).}
\skillrow{Idiomas}{Español nativo; inglés B2 (intermedio alto)}
\skillrow{Programación competitiva}{Participante en Google Code Jam (2019); equipo de selección de etapa de ACM-ICPC (2015--2017).}
}